\section{Выводы}

Выполнение данной лабораторной работы позволило на практике исследовать, как микрооптимизации влияют на реальную скорость работы структуры данных. Было видно, что даже при одинаковой асимптотической сложности конкретная реализация может заметно отличаться по времени выполнения из-за числа сравнений, рекурсивных вызовов и накладных расходов стандартных операций со строками.

Также были получены практические навыки работы с инструментами профилирования и анализа памяти. \texttt{gprof} помог определить наиболее дорогие участки кода, а \texttt{Valgrind} подтвердил отсутствие утечек и ошибок работы с динамической памятью. В результате программа стала быстрее, а её поведение -- более предсказуемым и корректным.

\pagebreak
